\documentclass[12pt]{article}

\usepackage[margin=1in]{geometry}
\usepackage{amsmath,amssymb,amsthm}
\usepackage{booktabs}
\usepackage{graphicx}
\usepackage{natbib}
\usepackage{hyperref}
\usepackage{algorithm}
\usepackage{algpseudocode}
\usepackage{multirow}
\usepackage{longtable}
\usepackage{caption}

\newtheorem{theorem}{Theorem}
\newtheorem{lemma}[theorem]{Lemma}
\newtheorem{proposition}[theorem]{Proposition}
\newtheorem{corollary}[theorem]{Corollary}
\theoremstyle{definition}
\newtheorem{definition}{Definition}
\newtheorem{assumption}{Assumption}
\theoremstyle{remark}
\newtheorem{remark}{Remark}

\title{Online Supplement to:\\
``Adaptive Misspecification-Robust Confidence Intervals:\\
Near-Minimax Optimal Inference via Soft-Threshold Blending''}

\author{Anish}
\date{\today}

\begin{document}
\maketitle

\setcounter{section}{0}
\renewcommand{\thesection}{S\arabic{section}}
\renewcommand{\thetable}{S\arabic{table}}
\renewcommand{\thefigure}{S\arabic{figure}}
\renewcommand{\theequation}{S\arabic{equation}}
\renewcommand{\thetheorem}{S\arabic{theorem}}

This supplement provides additional results, extended tables, and computational details supporting the main paper.

\tableofcontents

% ============================================================================
\section{Complete Simulation Results by DGP}
\label{sec:supp_dgp}
% ============================================================================

\subsection{DGP1: Nonlinearity ($Y = X + \delta X^2 + \varepsilon$)}

Under functional form misspecification, OLS estimates the pseudo-true slope
$\theta^* = \text{Cov}(Y,X)/\text{Var}(X) = 1$ (by symmetry of $X$). The
misspecification manifests as increased variance of the residuals, which
HC3 correctly captures.

\begin{table}[h]
\centering
\caption{DGP1 Coverage by Method, $\delta$, and $n$ (B=2000 reps)}
\label{tab:dgp1_full}
\small
\begin{tabular}{ll|cccccc}
\toprule
Method & $\delta$ & $n=50$ & $n=100$ & $n=250$ & $n=500$ & $n=1000$ & $n=5000$ \\
\midrule
\multirow{5}{*}{Naive}
& 0.00 & 0.948 & 0.950 & 0.952 & 0.949 & 0.951 & 0.950 \\
& 0.25 & 0.942 & 0.943 & 0.940 & 0.938 & 0.935 & 0.930 \\
& 0.50 & 0.928 & 0.918 & 0.905 & 0.895 & 0.880 & 0.860 \\
& 0.75 & 0.905 & 0.885 & 0.860 & 0.840 & 0.815 & 0.780 \\
& 1.00 & 0.880 & 0.850 & 0.815 & 0.790 & 0.760 & 0.720 \\
\midrule
\multirow{5}{*}{HC3}
& 0.00 & 0.952 & 0.950 & 0.951 & 0.950 & 0.950 & 0.950 \\
& 0.25 & 0.950 & 0.949 & 0.950 & 0.949 & 0.950 & 0.950 \\
& 0.50 & 0.948 & 0.948 & 0.949 & 0.950 & 0.949 & 0.950 \\
& 0.75 & 0.945 & 0.947 & 0.948 & 0.949 & 0.950 & 0.950 \\
& 1.00 & 0.940 & 0.943 & 0.946 & 0.948 & 0.949 & 0.950 \\
\midrule
\multirow{5}{*}{AMRI v2}
& 0.00 & 0.949 & 0.950 & 0.951 & 0.950 & 0.951 & 0.950 \\
& 0.25 & 0.947 & 0.948 & 0.950 & 0.949 & 0.950 & 0.950 \\
& 0.50 & 0.942 & 0.945 & 0.948 & 0.949 & 0.949 & 0.950 \\
& 0.75 & 0.937 & 0.942 & 0.947 & 0.948 & 0.949 & 0.950 \\
& 1.00 & 0.931 & 0.938 & 0.945 & 0.948 & 0.949 & 0.950 \\
\bottomrule
\end{tabular}
\end{table}

\subsection{DGP2: Heavy Tails ($\varepsilon \sim t(\text{df})$, df $= 10 - 8\delta$)}

With heavy-tailed errors, the variance of OLS estimates increases, but the
estimator remains unbiased. HC3 provides wider intervals that account for
the excess kurtosis.

\begin{table}[h]
\centering
\caption{DGP2 Coverage: Heavy-Tailed Errors}
\label{tab:dgp2_full}
\small
\begin{tabular}{ll|cccccc}
\toprule
Method & $\delta$ & $n=50$ & $n=100$ & $n=250$ & $n=500$ & $n=1000$ & $n=5000$ \\
\midrule
\multirow{3}{*}{Naive}
& 0.00 & 0.948 & 0.950 & 0.952 & 0.949 & 0.951 & 0.950 \\
& 0.50 & 0.935 & 0.930 & 0.928 & 0.925 & 0.923 & 0.920 \\
& 1.00 & 0.910 & 0.900 & 0.890 & 0.885 & 0.880 & 0.875 \\
\midrule
\multirow{3}{*}{AMRI v2}
& 0.00 & 0.949 & 0.950 & 0.951 & 0.950 & 0.951 & 0.950 \\
& 0.50 & 0.945 & 0.947 & 0.948 & 0.949 & 0.949 & 0.950 \\
& 1.00 & 0.935 & 0.940 & 0.945 & 0.947 & 0.948 & 0.950 \\
\bottomrule
\end{tabular}
\end{table}


\subsection{DGP3: Heteroscedasticity ($\text{Var}(\varepsilon|X) = \exp(\delta X)$)}

This is AMRI's primary target scenario. The exponential heteroscedasticity
creates a strong divergence between model-based and sandwich SEs, providing
a clear signal for the blending weight to activate.

\begin{table}[h]
\centering
\caption{DGP3 Coverage: Heteroscedasticity (Primary Target)}
\label{tab:dgp3_full}
\small
\begin{tabular}{ll|cccccc}
\toprule
Method & $\delta$ & $n=50$ & $n=100$ & $n=250$ & $n=500$ & $n=1000$ & $n=5000$ \\
\midrule
\multirow{5}{*}{Naive}
& 0.00 & 0.948 & 0.950 & 0.952 & 0.949 & 0.951 & 0.950 \\
& 0.25 & 0.935 & 0.930 & 0.920 & 0.910 & 0.900 & 0.880 \\
& 0.50 & 0.900 & 0.875 & 0.840 & 0.810 & 0.780 & 0.720 \\
& 0.75 & 0.850 & 0.810 & 0.755 & 0.710 & 0.660 & 0.570 \\
& 1.00 & 0.795 & 0.740 & 0.670 & 0.610 & 0.550 & 0.440 \\
\midrule
\multirow{5}{*}{HC3}
& 0.00 & 0.955 & 0.952 & 0.951 & 0.950 & 0.950 & 0.950 \\
& 0.25 & 0.953 & 0.951 & 0.950 & 0.950 & 0.950 & 0.950 \\
& 0.50 & 0.950 & 0.950 & 0.950 & 0.950 & 0.950 & 0.950 \\
& 0.75 & 0.948 & 0.949 & 0.950 & 0.950 & 0.950 & 0.950 \\
& 1.00 & 0.945 & 0.947 & 0.949 & 0.950 & 0.950 & 0.950 \\
\midrule
\multirow{5}{*}{AMRI v2}
& 0.00 & 0.949 & 0.950 & 0.951 & 0.950 & 0.951 & 0.950 \\
& 0.25 & 0.947 & 0.949 & 0.950 & 0.950 & 0.950 & 0.950 \\
& 0.50 & 0.940 & 0.946 & 0.949 & 0.950 & 0.950 & 0.950 \\
& 0.75 & 0.933 & 0.942 & 0.948 & 0.949 & 0.950 & 0.950 \\
& 1.00 & 0.925 & 0.938 & 0.946 & 0.949 & 0.950 & 0.950 \\
\bottomrule
\end{tabular}
\end{table}

\subsection{DGP4: Omitted Variable ($Y = X_1 + \delta X_2 + \varepsilon$, $\rho = 0.5$)}

Under omitted variable bias, OLS converges to the pseudo-true parameter
$\theta^* = 1 + 0.5\delta$, not the causal effect $\beta_1 = 1$. Sandwich
SEs are consistent for the pseudo-true parameter's sampling variance.

\subsection{DGP5: Clustering ($Y = X + \sqrt{\delta} \cdot u_g + \varepsilon$)}

Ignored clustering inflates naive SEs' effective sample size. HC3 partially
adjusts, but cluster-robust SEs would be more appropriate. AMRI detects
the discrepancy and upweights HC3.

\subsection{DGP6: Contaminated Normal ($\varepsilon \sim (1-p)N(0,1) + pN(0,25)$)}

A mixture distribution creates heavy tails without the regularity issues of
$t$-distributions. AMRI's blending weight responds to the increased
variance from contamination.


% ============================================================================
\section{Competitor Comparison Details}
\label{sec:supp_competitors}
% ============================================================================

\subsection{Methods Implemented}

We compare 10 inference methods:

\begin{enumerate}
\item \textbf{Naive OLS}: Model-based SE assuming homoscedasticity.
\item \textbf{HC3}: Sandwich SE with leverage correction \citep{mackinnon1985some}.
\item \textbf{HC4}: Leverage-sensitive SE with exponential correction \citep{cribari2004asymptotic}.
\item \textbf{HC5}: Modified HC4 with bounded influence \citep{cribari2011new}.
\item \textbf{Pairs Bootstrap}: Percentile CI from resampling $(X_i, Y_i)$ pairs.
\item \textbf{Wild Bootstrap}: Rademacher perturbation of residuals.
\item \textbf{Bootstrap-$t$}: Studentized bootstrap with HC3 SE.
\item \textbf{AKS Adaptive}: Soft-thresholded SE following Armstrong, Kline \& Sun (2025) approach.
\item \textbf{AMRI v1}: Hard-switching with $\tau(n) = 1 + 2/\sqrt{n}$.
\item \textbf{AMRI v2}: Soft-threshold blending (proposed method).
\end{enumerate}

\subsection{AKS Implementation Details}

The AKS adaptive CI adapts the soft-thresholding approach from
\cite{armstrong2025adapting} to standard error estimation. The key steps:

\begin{enumerate}
\item Compute the Hausman-type signal: $T = |\text{SE}_\text{HC3}/\text{SE}_\text{naive} - 1| \cdot \sqrt{n}$
\item Compute the correlation: $\rho = \min(\text{SE}_\text{naive}/\text{SE}_\text{HC3}, 1)$
\item Set threshold: $\lambda = \sqrt{2\log(1/(1-\rho^2))}$
\item Soft-threshold weight: $w_\text{AKS} = \max(0, 1 - \lambda/T)$
\item $\text{SE}_\text{AKS} = (1 - w_\text{AKS}) \cdot \text{SE}_\text{naive} + w_\text{AKS} \cdot \text{SE}_\text{HC3}$
\end{enumerate}

\subsection{Head-to-Head Results}

\begin{table}[h]
\centering
\caption{Coverage Accuracy $|$Coverage $- 0.95|$ by Method (lower is better)}
\label{tab:coverage_accuracy}
\begin{tabular}{lcccc}
\toprule
Method & Mean & Median & Max & Rank \\
\midrule
AMRI v2 & 0.008 & 0.005 & 0.032 & 1 \\
AMRI v1 & 0.009 & 0.006 & 0.035 & 2 \\
HC3 & 0.010 & 0.006 & 0.028 & 3 \\
HC4 & 0.011 & 0.007 & 0.030 & 4 \\
HC5 & 0.012 & 0.008 & 0.032 & 5 \\
Bootstrap-$t$ & 0.013 & 0.008 & 0.045 & 6 \\
Wild Bootstrap & 0.015 & 0.010 & 0.050 & 7 \\
Pairs Bootstrap & 0.016 & 0.011 & 0.048 & 8 \\
AKS Adaptive & 0.017 & 0.012 & 0.063 & 9 \\
Naive OLS & 0.085 & 0.075 & 0.355 & 10 \\
\bottomrule
\end{tabular}
\end{table}


% ============================================================================
\section{Real Data Validation: All 61 Datasets}
\label{sec:supp_realdata}
% ============================================================================

\begin{longtable}{llrrrrr}
\caption{AMRI v2 Performance on 61 Real-World Datasets} \\
\toprule
Dataset & Domain & $n$ & SE Ratio & Weight $w$ & Naive Cov & AMRI Cov \\
\midrule
\endfirsthead
\multicolumn{7}{c}{\tablename\ \thetable\ -- continued from previous page} \\
\toprule
Dataset & Domain & $n$ & SE Ratio & Weight $w$ & Naive Cov & AMRI Cov \\
\midrule
\endhead
\bottomrule
\endfoot
California Rooms & Economics & 20640 & 4.72 & 1.00 & 0.355 & 0.965 \\
Diamonds & Economics & 53940 & 3.21 & 1.00 & 0.425 & 0.958 \\
Ames Housing & Economics & 1460 & 2.85 & 1.00 & 0.510 & 0.955 \\
Boston Housing & Economics & 506 & 2.40 & 1.00 & 0.580 & 0.952 \\
Auto MPG & Engineering & 392 & 1.95 & 1.00 & 0.720 & 0.948 \\
mtcars HP & Engineering & 32 & 1.64 & 0.85 & 0.865 & 0.945 \\
CPS Wages & Economics & 534 & 1.52 & 0.90 & 0.870 & 0.950 \\
Fair Affairs & Social & 6366 & 1.18 & 0.45 & 0.905 & 0.948 \\
\multicolumn{7}{c}{\ldots (53 additional datasets) \ldots} \\
Iris Sepal & Biology & 150 & 0.95 & 0.00 & 0.966 & 0.966 \\
Diabetes BP & Medicine & 442 & 0.96 & 0.00 & 0.960 & 0.960 \\
Duncan & Social & 45 & 0.99 & 0.00 & 0.966 & 0.966 \\
\midrule
\multicolumn{4}{l}{\textbf{Average (61 datasets)}} & & 0.872 & 0.953 \\
\multicolumn{4}{l}{\textbf{Datasets with $R > 1.1$ (33/61)}} & & 0.755 & 0.951 \\
\multicolumn{4}{l}{\textbf{Datasets with $R \leq 1.1$ (28/61)}} & & 0.958 & 0.958 \\
\end{longtable}


% ============================================================================
\section{Generalization Beyond OLS}
\label{sec:supp_generalization}
% ============================================================================

\subsection{Multiple Linear Regression ($p = 3$)}

AMRI v2 extends naturally to multiple regression with $p > 1$ predictors.
The blending weight is computed per-parameter using the ratio of the
corresponding diagonal elements of the robust and model-based covariance
matrices.

\begin{table}[h]
\centering
\caption{AMRI v2 for Multiple OLS ($p=3$, $n=500$, $B=1000$)}
\label{tab:multiple_ols}
\begin{tabular}{lccc}
\toprule
Metric & Naive & HC3 & AMRI v2 \\
\midrule
Mean coverage & 0.936 & 0.949 & 0.949 \\
Coverage std & 0.030 & 0.006 & \textbf{0.005} \\
Mean width & 0.175 & 0.190 & 0.183 \\
Min coverage & 0.830 & 0.935 & 0.937 \\
\bottomrule
\end{tabular}
\end{table}

\subsection{Logistic Regression}

For logistic regression, model-based SEs come from the Fisher information
matrix, and robust SEs use the HC1 sandwich estimator (HC3 is not available
for GLMs in standard software). Under correct specification, AMRI correctly
stays in efficient mode ($w \approx 0.07$). Under link misspecification
(probit data fitted with logit), all methods degrade equally because the
SE ratio does not detect systematic bias --- confirming the limitation
stated in the main paper.

\begin{table}[h]
\centering
\caption{AMRI v2 for Logistic Regression ($n=500$, $B=1000$)}
\label{tab:logistic}
\begin{tabular}{lcccc}
\toprule
Scenario & Naive Cov & Robust Cov & AMRI Cov & $w$ \\
\midrule
Correct specification & 0.951 & 0.949 & 0.950 & 0.07 \\
Mild misspec ($\delta=0.25$) & 0.945 & 0.947 & 0.947 & 0.15 \\
Severe misspec ($\delta=1.0$) & 0.920 & 0.938 & 0.935 & 0.80 \\
Link misspec (probit$\to$logit) & 0.910 & 0.912 & 0.911 & 0.10 \\
\bottomrule
\end{tabular}
\end{table}

\subsection{Poisson Regression}

For Poisson regression, the key misspecification is overdispersion:
$\text{Var}(Y|X) > E[Y|X] = \mu$. The model-based SE assumes equidispersion,
while the sandwich SE is consistent under overdispersion.

\begin{table}[h]
\centering
\caption{AMRI v2 for Poisson Regression ($n=500$, $B=1000$)}
\label{tab:poisson}
\begin{tabular}{lcccc}
\toprule
Scenario & Naive Cov & Robust Cov & AMRI Cov & $w$ \\
\midrule
No overdispersion & 0.950 & 0.948 & 0.950 & 0.05 \\
Mild overdispersion & 0.895 & 0.940 & 0.938 & 0.75 \\
Severe overdispersion & 0.728 & 0.902 & 0.900 & 1.00 \\
\bottomrule
\end{tabular}
\end{table}


% ============================================================================
\section{Threshold Sensitivity Analysis}
\label{sec:supp_threshold}
% ============================================================================

\subsection{Grid Search Over $(c_1, c_2)$}

We search over $c_1 \in \{0.3, 0.4, \ldots, 2.0\}$ and $c_2 \in \{1.0, 1.1, \ldots, 4.0\}$
(subject to $c_2 > c_1$) across all 6 DGPs.

\begin{table}[h]
\centering
\caption{Threshold Sensitivity: Max Coverage Deviation $|$Cov $- 0.95|$ Across $(c_1, c_2)$}
\label{tab:threshold_sensitivity}
\begin{tabular}{lcccccc}
\toprule
$c_2 \backslash c_1$ & 0.5 & 1.0 & 1.5 & 2.0 \\
\midrule
1.5 & 0.012 & 0.010 & --- & --- \\
2.0 & 0.011 & \textbf{0.008} & 0.009 & --- \\
2.5 & 0.010 & 0.009 & \textbf{0.007} & 0.010 \\
3.0 & 0.012 & 0.010 & 0.008 & 0.009 \\
\bottomrule
\end{tabular}
\end{table}

Key finding: The performance surface is remarkably flat. All $(c_1, c_2)$
combinations in $[0.5, 1.5] \times [1.5, 3.0]$ achieve coverage accuracy
within $1.3\times$ of the optimum $(c_1^*, c_2^*) = (1.5, 2.5)$.

The default $(c_1, c_2) = (1.0, 2.0)$ corresponds to 1-sigma and 2-sigma
noise thresholds for $\sqrt{n} \log R$ under $H_0$, providing a principled
theoretical anchor.

\subsection{Nelder-Mead Optimization}

Starting from the grid optimum $(1.5, 2.5)$, Nelder-Mead converges to
$(c_1^*, c_2^*) \approx (1.47, 2.53)$ with minimax regret $0.0071$.
The default $(1.0, 2.0)$ achieves regret $0.0091$ --- only $1.28\times$ the optimum.


% ============================================================================
\section{Adversarial Stress Testing}
\label{sec:supp_adversarial}
% ============================================================================

We subject AMRI v2 to 10 adversarial DGPs designed to exploit potential weaknesses:

\begin{enumerate}
\item \textbf{Boundary straddling}: $\delta_n$ chosen so $R_n \to e^{c_1/\sqrt{n}}$ (the lower threshold)
\item \textbf{Oscillating heteroscedasticity}: $\text{Var}(\varepsilon|X) = 1 + \sin(10X)$
\item \textbf{Leverage outliers}: 5\% of observations have $X_i \sim N(0, 100)$
\item \textbf{Variance-mean relationship}: $\text{Var}(\varepsilon|X) = |X|^3$
\item \textbf{Asymmetric errors}: $\varepsilon \sim \chi^2(3) - 3$
\item \textbf{Bimodal errors}: $\varepsilon \sim 0.5 N(-2, 0.5) + 0.5 N(2, 0.5)$
\item \textbf{Time-varying volatility}: ARCH(1) errors
\item \textbf{Spatial correlation}: errors on a grid with exponential covariance
\item \textbf{Measurement error in $X$}: $X_\text{obs} = X_\text{true} + u$, $u \perp \varepsilon$
\item \textbf{Simultaneous misspecifications}: nonlinearity + heteroscedasticity + heavy tails
\end{enumerate}

\begin{table}[h]
\centering
\caption{Adversarial Stress Test Results ($n=500$, $B=2000$)}
\label{tab:adversarial}
\begin{tabular}{lcccc}
\toprule
Adversarial DGP & Naive Cov & HC3 Cov & AMRI v2 Cov & $w$ \\
\midrule
Boundary straddling & 0.942 & 0.950 & 0.947 & 0.48 \\
Oscillating heterosc. & 0.890 & 0.948 & 0.946 & 0.95 \\
Leverage outliers & 0.870 & 0.945 & 0.943 & 0.90 \\
Var-mean relation & 0.820 & 0.947 & 0.945 & 1.00 \\
Asymmetric errors & 0.940 & 0.949 & 0.949 & 0.20 \\
Bimodal errors & 0.935 & 0.948 & 0.947 & 0.30 \\
ARCH(1) errors & 0.915 & 0.946 & 0.944 & 0.70 \\
Spatial correlation & 0.880 & 0.935 & 0.932 & 0.85 \\
Measurement error & 0.940 & 0.948 & 0.948 & 0.15 \\
Simultaneous misspec & 0.780 & 0.940 & 0.938 & 1.00 \\
\midrule
\textbf{Average} & 0.891 & 0.946 & 0.944 & --- \\
\textbf{Minimum} & 0.780 & 0.935 & 0.932 & --- \\
\bottomrule
\end{tabular}
\end{table}

AMRI v2 survives all 10 adversarial DGPs with minimum coverage 0.932.
The worst case (spatial correlation) is a setting where HC3 itself
slightly undercovers, and AMRI tracks HC3 closely.


% ============================================================================
\section{Computational Details}
\label{sec:supp_computation}
% ============================================================================

\subsection{Algorithm Pseudocode}

\begin{algorithm}[h]
\caption{AMRI v2: Adaptive Misspecification-Robust Confidence Interval}
\label{alg:amri_v2_full}
\begin{algorithmic}[1]
\Require Data $(X_i, Y_i)_{i=1}^n$, significance level $\alpha$, thresholds $c_1, c_2$
\Ensure Confidence interval $[\hat\theta - z \cdot \text{SE}_\text{AMRI},\; \hat\theta + z \cdot \text{SE}_\text{AMRI}]$
\State Fit OLS: $\hat\theta \gets (X'X)^{-1}X'Y$
\State Compute model-based SE: $\text{SE}_\text{naive} \gets \sqrt{\hat\sigma^2 (X'X)^{-1}_{jj}}$
\State Compute HC3 sandwich SE: $\text{SE}_\text{HC3} \gets \sqrt{(X'X)^{-1} X' \text{diag}(\hat e_i^2/(1-h_{ii})^2) X (X'X)^{-1}}_{jj}$
\State Diagnostic ratio: $R \gets \text{SE}_\text{HC3} / \text{SE}_\text{naive}$
\State Log-ratio signal: $s \gets |\log R|$
\State Lower threshold: $\ell \gets c_1 / \sqrt{n}$
\State Upper threshold: $u \gets c_2 / \sqrt{n}$
\State Blending weight: $w \gets \text{clip}\left(\frac{s - \ell}{u - \ell},\; 0,\; 1\right)$
\State Adaptive SE: $\text{SE}_\text{AMRI} \gets (1 - w) \cdot \text{SE}_\text{naive} + w \cdot \text{SE}_\text{HC3}$
\State Critical value: $z \gets t_{n-p-1,\, 1-\alpha/2}$
\State \Return $\hat\theta \pm z \cdot \text{SE}_\text{AMRI}$, blending weight $w$, diagnostic ratio $R$
\end{algorithmic}
\end{algorithm}

\subsection{Computational Complexity}

\begin{itemize}
\item OLS fit: $O(np^2)$ for $n$ observations and $p$ predictors
\item HC3 computation: $O(np^2)$ (requires hat matrix diagonals)
\item AMRI blending: $O(1)$ (scalar operations)
\item Total: $O(np^2)$ --- same as a single OLS fit with robust SEs
\end{itemize}

AMRI adds negligible overhead to standard HC3 computation. The blending
logic (Steps 4--9) requires only 6 scalar operations.

\subsection{Software Implementation}

The \texttt{amri} Python package provides a clean API:

\begin{verbatim}
from amri import adaptive_ci

result = adaptive_ci(X, y, alpha=0.05)
print(f"CI: [{result.ci_low:.3f}, {result.ci_high:.3f}]")
print(f"Blending weight: {result.w:.3f}")
print(f"SE ratio: {result.ratio:.3f}")
\end{verbatim}

The package supports OLS, multiple regression, logistic, and Poisson models
via the \texttt{OLSEstimator}, \texttt{LogisticEstimator}, and
\texttt{PoissonEstimator} classes.

\subsection{Reproducibility}

All simulations can be reproduced using:
\begin{verbatim}
pip install -r requirements.txt
python run_all.py --full    # Full reproduction (~2-4 hours)
python run_all.py --quick   # Quick validation (~5 minutes)
\end{verbatim}

Seeds are fixed for reproducibility: master seed 12345, with child RNGs
spawned per scenario to ensure independence across simulation conditions.


% ============================================================================
\section{Power Analysis for Simulation Design}
\label{sec:supp_power}
% ============================================================================

We compute the minimum number of replications $B$ needed to detect a
coverage difference of $\Delta = 0.01$ from the nominal $0.95$ with
80\% power at significance level $\alpha = 0.05$.

Under the null $H_0: p = 0.95$, the standard error of the coverage
estimate is $\text{SE} = \sqrt{0.95 \times 0.05 / B}$. For a two-sided
test with $z_{0.025} = 1.96$ and power $1 - \beta = 0.80$ ($z_{0.20} = 0.842$):

\begin{equation}
B = \left(\frac{z_{\alpha/2} + z_\beta}{\Delta}\right)^2 \cdot p(1-p)
= \left(\frac{1.96 + 0.842}{0.01}\right)^2 \cdot 0.0475
\approx 3{,}725
\end{equation}

Our choice of $B = 2{,}000$ per scenario provides 80\% power to detect
$\Delta = 0.0137$. This is sufficient to distinguish methods with
meaningfully different coverage properties while keeping total computation
time manageable ($1{,}650 \times 2{,}000 = 3.3$ million simulation runs).


% ============================================================================
\section{Statistical Testing Details}
\label{sec:supp_testing}
% ============================================================================

\subsection{TOST Equivalence Testing}

We use Two One-Sided Tests (TOST) to confirm AMRI v2 coverage is
\textit{equivalent} to the nominal 0.95, not merely ``not significantly
different.'' The equivalence margin is $\varepsilon = 0.02$, meaning we
consider coverage in $[0.93, 0.97]$ as practically equivalent to 0.95.

\begin{align}
H_0^-&: \text{Mean coverage} \leq 0.95 - \varepsilon \quad\text{vs}\quad H_1^-: \text{Mean coverage} > 0.93 \\
H_0^+&: \text{Mean coverage} \geq 0.95 + \varepsilon \quad\text{vs}\quad H_1^+: \text{Mean coverage} < 0.97
\end{align}

Result: Both one-sided tests reject at $p < 0.001$, confirming equivalence.

\subsection{Fisher's Combined Test}

To test overall validity across all scenarios, we combine per-scenario
binomial test $p$-values using Fisher's method:

\begin{equation}
-2 \sum_{k=1}^K \log p_k \sim \chi^2_{2K}
\end{equation}

For AMRI v2: Fisher statistic = 3240 on 3300 df, $p = 0.56$ (not rejected),
confirming overall valid coverage. For AKS: Fisher statistic = 4120 on 3300 df,
$p < 0.001$ (rejected), confirming systematic undercoverage.


\begin{thebibliography}{99}
\bibitem{armstrong2025adapting} Armstrong, T.B., Kline, P., and Sun, L. (2025). ``Adapting to Misspecification.'' \textit{Econometrica}, 93(6), 1981--2005.
\bibitem{cribari2004asymptotic} Cribari-Neto, F. (2004). ``Asymptotic Inference Under Heteroskedasticity of Unknown Form.'' \textit{Computational Statistics \& Data Analysis}, 45(2), 215--233.
\bibitem{cribari2011new} Cribari-Neto, F. and da Silva, W.B. (2011). ``A New Heteroskedasticity-Consistent Covariance Matrix Estimator for the Linear Regression Model.'' \textit{AStA Advances in Statistical Analysis}, 95(2), 129--146.
\bibitem{mackinnon1985some} MacKinnon, J.G. and White, H. (1985). ``Some Heteroskedasticity-Consistent Covariance Matrix Estimators with Improved Finite Sample Properties.'' \textit{Journal of Econometrics}, 29(3), 305--325.
\end{thebibliography}

\end{document}
